%%%%%%%%%%%%%%%%%%%%%%%%%%%%%%%%%%%%%%%%%%%%%%%%
% Gerätenutzung: was die farbigen Maschinenschilder bedeuten (Aushang A3)
% CC-BY-SA 3.0, FAU FabLab
% Früher Hinweisschilder_divers/geraetenutzung.pdf (Illustrator), jetzt LaTeX.
% Die Schilder selbst: github.com/fau-fablab/hinweisschilder
%%%%%%%%%%%%%%%%%%%%%%%%%%%%%%%%%%%%%%%%%%%%%%%%
\documentclass[a3paper]{scrartcl}
\usepackage{fablab-aushang}
\Fabcube

% Beispielschild:  \Beispiel{Farbe}{Gerät}{Regel}{Kleingedrucktes}
% (Breite fest, weil minipage \textwidth und \linewidth ändert)
\newlength{\BspBreite}
\AtBeginDocument{\setlength{\BspBreite}{0.3\textwidth}}
\newcommand{\Beispiel}[4]{%
  \tikz\node[fill=#1, draw=black!70, line width=1pt, rounded corners=0.25\SchildText,
    text width=\BspBreite, align=center, inner sep=0.4\SchildText]
    {{\bfseries\fontsize{1.05\SchildText}{1.2\SchildText}\selectfont #2\par}
     \vspace{0.3\SchildText}
     {\fontsize{0.95\SchildText}{1.1\SchildText}\selectfont #3\par}
     \vspace{0.3\SchildText}
     {\fontsize{0.5\SchildText}{0.6\SchildText}\selectfont #4\par}};}
% eine Zeile: Beispielschild links, Regeln rechts
\newcommand{\Stufe}[5]{%
  \par\noindent
  \begin{minipage}[c]{0.37\linewidth}\Beispiel{#1}{#2}{#3}{#4}\end{minipage}\hfill
  \begin{minipage}[c]{0.6\linewidth}\begin{itemize}#5\end{itemize}\end{minipage}\par
  \vspace{0.4\SchildText}{\color{black!40}\rule{\linewidth}{1pt}}\par}

\begin{document}
\Schild[0.66]
\SchildKopf[0.85]{Gerätenutzung}

Alle Geräte und Maschinen sind folgendermaßen farbig beschildert:

\Stufe{ampelGruen}{Staubsauger\\ BLITZ-BLANK42}{}{Zum Nass- und Trockensaugen geeignet.}
  {\item Darf \textbf{ohne} Einweisung eigenständig benutzt werden
   \item Bei Unklarheiten nachfragen}
\Stufe{ampelGelb}{Ersatzluftblase}
  {Selbstständige Benutzung nur mit unterschriebener Wasserwaagen-Einweisung}
  {Bei Interesse frage einfach einen Betreuer!}
  {\item Darf \textbf{nach unterschriebener Einweisung} eigenständig benutzt werden
   \item Einweisung immer beachten}
\Stufe{ampelRot}{CNC-Pinsel\\ Norbert iPaint}{Benutzung nur nach Rücksprache.}
  {Im Moment gibt es leider noch keine Einweisung zu diesem Gerät.
   Bitte sprich mit einem Betreuer.}
  {\item Darf \textbf{nur nach Rücksprache} mit Betreuern benutzt werden
   \item Besondere Regelungen sind zu beachten}

\vspace{0.4\SchildText}
Für alle Nutzer und Gegenstände gilt jederzeit:
\begin{itemize}
  \item Maschinenbeschilderung beachten
  \item FabLab-Eigentum pfleglich behandeln und Beschädigungen einem Betreuer melden
  \item Anweisungen der Betreuer befolgen
  \item Menschenverstand einsetzen
  \item Arbeitsplätze und Geräte nach der Nutzung sauber und ordentlich hinterlassen
\end{itemize}

\end{document}
